%
% 09_schlusswort.tex -- Schlusswort
%
% (c) 2026 Dominik Castelberg, Ostschweizer Fachhochschule
%
% !TEX root = ../../buch.tex
% !TEX encoding = UTF-8
%
\section{Der Satz von Jordan: Ein Fazit
\label{jordan:section:schlusswort}}
\kopfrechts{Der Satz von Jordan: Ein Fazit}

Was mit einer einfachen Frage, welche wir alle intuitiv korrekt beantworten können, begann, hat sich zu einer erstaunlich tiefen Erkundung der Geometrie, Algebra und Topologie entwickelt.
Über die Alexander-Dualität hat die Fragestellung, die der Satz von Jordan aufwirft, weitreichende Konsequenzen in der Algebra und Topologie, welche wir in diesem Paper nur oberflächlich berührt haben, hervorgebracht.

Wie in den ersten Sätzen dieses Papers erwähnt, sind manche oberflächlich trivial wirkende Sätze in der Mathematik durchaus einen Beweis wert.
Und einige wenige dieser Sätze, wie der Satz von Jordan, sind sogar so tiefgründig, dass ihre Beweise im Allgemeinen ganze Disziplinen der Mathematik revolutionieren.